



\pdfbookmark[0]{PORTADA}{prt} 
\thispagestyle{empty}

\begin{figure}[h]
\setbox0 \vbox{\begin{flushleft}
\vspace*{-18pt}\hspace{-10mm}\includegraphics[width = 2.5cm]{escudo-eps-converted-to.pdf}\\
\end{flushleft}}
\wd0 = 0pt \ht0 = 0pt \box0
\end{figure} \hangindent= 18mm \hangafter=-3\vspace{-10mm}
\noindent Universidad de San Carlos de Guatemala\\[5pt]
Escuela de Ciencias Físicas y Matemáticas \\[5pt]
Departamento de Física

\vfill

\begin{center}
{\setlength{\baselineskip}{1.3\baselineskip} \large \textbf{\MakeUppercase{Protocolo de tesis: Operadores de Kraus para mediciones difusas en sistemas de \boldmath{$n$} partículas}}\par}

\vfill

\textbf{Rubí Esmeralda Ramírez Milián}\\[6pt]
Asesorado por: \\
\textit{Dr.\@ Carlos Francisco Pineda Zorrilla e Ing. Rodolfo Samayoa }  

\vfill

Guatemala, {noviembre 2023}
\end{center}

%%%%%%%%%%%%PAGINA EN BLANCO%%%%%%%%%%%%%%%%%%%%%%%%%%%%

\newpage

\textcolor[rgb]{1,1,1}{.} \thispagestyle{empty}
\begin{comment}
%%%%%%%%%%%%%% SEGUNDA PORTADA %%%%%%%%%%%%%%%%%%%%%%%%%%
\newpage

\pdfbookmark[1]{IDENTIFICACIÓN}{idn}
\thispagestyle{empty}

\begin{center}
UNIVERSIDAD DE SAN CARLOS DE GUATEMALA
\end{center}

\begin{figure}[h]
  \begin{center}
    \includegraphics[width=4.15cm]{escudo-eps-converted-to.pdf}\\
  \end{center}
\end{figure}

\begin{center}
\vspace{-0.5mm}ESCUELA DE CIENCIAS FÍSICAS Y MATEMÁTICAS

\vfill

\recuadro{\MakeUppercase{Operadores de Kraus para mediciones difusas en sistemas de $N$ partículas}}
\vfill

TRABAJO DE GRADUACIÓN \\[3pt] PRESENTADO A LA JEFATURA DEL \\[3pt]
DEPARTAMENTO DE FÍSICA\\[3pt] POR \\[1.15cm]
\textbf{\MakeUppercase{Rubí Esmeralda Ramírez Milián}} \\[6pt]
ASESORADO POR {\MakeUppercase{Carlos Francisco Pineda Zorrilla}} 

\vfill
AL CONFERÍRSELE EL TÍTULO DE \\[6pt]
\textbf{LICENCIADA EN FÍSICA APLICADA}

\vfill

GUATEMALA, {\MakeUppercase{2023}}

\end{center}


%%%%%%%%%PAGINA VACÍA%%%%%%%%%%%%%%%%%%%%%%%%%%%%%%%%%%%
\newpage
\textcolor[rgb]{1,1,1}{.} 

\thispagestyle{empty}




%%%%%%%%%%%%%%% PAGINA DEL PRIVADO%%%%%%%%%%%%%%%%%%%%%%%%%%%%%
\newpage

\pdfbookmark[1]{NÓMINA DEL CONSEJO DIRECTIVO}{ndj}
\thispagestyle{empty}

\begin{center}
UNIVERSIDAD DE SAN CARLOS DE GUATEMALA\\[6pt] ESCUELA DE CIENCIAS FÍSICAS Y MATEMÁTICAS
\end{center}

\begin{figure}[h]
  \begin{center}
    \includegraphics[width=4.15cm]{escudo-eps-converted-to.pdf}\\
  \end{center}
\end{figure}

\begin{center}
\textbf{CONSEJO DIRECTIVO}
\end{center}

{\doublespacing{}
\begin{tabbing}
  SECRETARIO ACADÉMICO xs \= xxxxxxxxxxxxxxxxxxxxxxxxxxxxxxxxxxxx  \kill
  DIRECTOR \> Marcelo Ixquiac \\
%  VOCAL I \>  vocal-1\\
%  VOCAL II \> vocal-2 \\
%  VOCAL III \> vocal-3 \\
%  VOCAL IV \> vocal-4 \\
%  VOCAL V \> vocal-5 \\
  SECRETARIO ACADÉMICO \> Edgar Cifuentes
\end{tabbing}

\vspace{24pt}

\begin{center}
\textbf{TRIBUNAL QUE PRACTICÓ EL EXAMEN GENERAL PRIVADO}
\end{center}

\begin{tabbing}
  SECRETARIO ACADÉMICOxs \= xxxxxxxxxxxxxxxxxxxxxxxxxxxxxxxxxxxx  \kill
%  DIRECTOR \> decano en funciones al hacer privado \\
  EXAMINADOR \> 1 \\
  EXAMINADOR \> 2 \\
  EXAMINADOR \> 3 \\
%  SECRETARIO ACADÉMICO \> secretario en funciones al hacer privado
\end{tabbing}
\par}

\newpage
\textcolor[rgb]{1,1,1}{.} 

\thispagestyle{empty}
\newpage









%%%%%%%%%%%%%%%%%%%%%%%%%%%%%%%%%%%%%%%%%%%%%%%%%%%%%%%%%%%%%%%%%%%%%%%%%%%%%%%%%%%%%%%%%%%%%%%
\begin{titlepage}
\vbox{ }
\vbox{ }
\begin{center}
	% Title
	\HRule{} \\[0.4cm]
	{ \huge \bfseries Operadores de Kraus en mediciones difusas para dos partículas}\\[0.4cm]
	\HRule{} \\[1.5cm]
	
	
	\Large
	%\emph{Autora}\\
	\textit{Rubí Esmeralda Ramírez Milián}\\[1.5cm]
% Upper part of the page
%\includegraphics[width=0.50\textwidth]{Images/logos.jpeg}\\
%\includegraphics[width=0.30\textwidth]{Images/logocolor.png}\\[1cm]
\textsc{\Large Universidad de San Carlos de Guatemala}\\[0.2cm]
\textsc{\Large Escuela de Ciencias Físicas y Matemáticas}\\[1.5cm]
\textsc{\large Informe final de año de prácticas}\\[1.5cm]
\vbox{ }


% Author
\large
Supervisado por:\\
\emph{Dr.\ Carlos Pineda (IF-UNAM) e Ing.\ Rodolfo Samayoa (ECFM-USAC)}\\
\vfill

% Bottom of the page
{\large\emph{18 de marzo 2023}}
\end{center}
\end{titlepage}
\end{comment}